\documentclass[12pt]{article}
\usepackage[shortlabels]{enumitem}
\usepackage{float}
\usepackage{xcolor}
\usepackage[a4paper, top=1.5cm, bottom=1.5cm, left=2cm, right=2cm]{geometry}
\usepackage{parskip}
\usepackage{setspace}
\usepackage{lmodern}
\usepackage[T1]{fontenc}
\usepackage[brazil]{babel}
\usepackage{hyperref}
\usepackage{graphicx}
\usepackage{amsmath}
\usepackage{amssymb}
\usepackage{amsfonts}
\usepackage{dsfont}
\usepackage{float}
\title{Lista 6 - Estatística para Administração - 2025.2}
\begin{document}
\date{}
\maketitle

1) Qual a diferença entre variáveis aleatórias discretas e contínuas? De um exemplo de variável aleatória discreta e outro exemplo de uma contínua. 

\vspace{3px}

2) Quantos grupos de 4 alunos podem ser formados a partir de uma turma de 30 alunos?

\vspace{3px}

3) Uma empresa está testando uma nova campanha de marketing digital voltada para fortalecer a marca empregadora (Employer Branding).
 Em cada visita ao site da empresa, existe $25\%$ de chance do visitante clicar no botão \textit{“Trabalhe Conosco”}. 

    Durante uma semana, $18$ visitantes chegam ao site por meio da nova campanha.

    \begin{enumerate}[a)]
        \item Modele a variável aleatória $X$ definida como o número de visitantes que clicam no botão \textit{“Trabalhe Conosco”}. Qual a distribuição de $X$? Por que?
        \item Calcule a probabilidade de exatamente $5$ dos $18$ visitantes clicarem no botão.
        \item Qual a probabilidade de nenhum cliente clicar no botão?
        \item Calcule a probabilidade de pelo menos $2$ dos $18$ visitantes clicarem no botão.
    \end{enumerate}

\vspace{3px}

4) Seja $X_1, X_2, X_3, X_4$ variáveis aleatórias I.I.D tal que $X_1 \sim Bernoulli(0,3)$. Qual a distribuição de $Y=\displaystyle\sum_{i=1}^4 X_i$ ?

\vspace{3px}

5) Seja $Z$ uma variável aleatória tal que $Z\sim N(0,1)$, calcule:

\begin{enumerate}[a)]
    \item $\mathds{P}(Z>1)$
    \item $\mathds{P}(-1 \leq Z<1)$
    \item $\mathds{P}(Z>2| Z> 1)$
    \item $\mathds{P}(Z\leq 1,64)$
    \item $\mathds{P}(Z \leq 2,56)$
    \item $\mathds{P}(-2,56 \leq Z \leq 2,56)$
    \item $\mathds{P}(Z>2| Z< -1)$
    \item $\mathds{P}(Z>2| Z< 3)$
    \item $\mathds{P}(Z<-2)$
    \item $\mathds{P}(-4<Z<4)$
    \item $\mathds{P}(Z>12)$
\end{enumerate}

\vspace{3px}

6) Suponha que a altura dos homens de 25 anos de idade, em cm, seja uma variável aleatória normal com parâmetros $\mu = 180$ e $\sigma^2=16$. 
Qual a probabilidade de um homem de 25 anos de idade ter mais de 1,88 de altura?

\vspace{3px}

7) Se X é uma variável aleatória com distribuição Normal com parâmetros $\mu=10$ e $\sigma^2=36$, calcule:

\begin{enumerate}[a)]
    \item $\mathds{P}(X>5)$
    \item $\mathds{P}(4<X<16)$
    \item $\mathds{P}(X<8)$
    \item $\mathds{P}(X<20)$
    \item $\mathds{P}(X>16)$
    \item $\mathds{P}(X>5 | X> 2)$
    \item $x$ tal que $\mathds{P}(X>x) = 0.75$
    \item $\mathds{P}(|X-10| \leq 5)$
\end{enumerate}

\vspace{3px}


8) 
A área de Recursos Humanos de uma grande empresa do setor varejista realiza anualmente uma pesquisa de satisfação salarial
 com seus colaboradores. Estudos setoriais anteriores indicam que o desvio padrão populacional para a insatisfação salarial
(escala de 0 a 100) é de 12 pontos.

Em uma amostra aleatória de 64 colaboradores da empresa, 
a insatisfação salarial média foi de 58 pontos em uma escala onde valores mais altos indicam maior insatisfação.

\begin{enumerate}[a)]
    \item Construa um intervalo de confiança de 90\% para a verdadeira insatisfação salarial média populacional dos colaboradores desta empresa. Interprete o intervalo de confiança encontrado.
    \item Construa um intervalo de confiança de 95\% para a verdadeira insatisfação salarial média populacional dos colaboradores desta empresa. Interprete o intervalo de confiança encontrado.
    \item Discuta a influência do tamanho amostral e do nível de confiança na amplitude (comrpimento) do intervalo. 
\end{enumerate}

\vspace{3px}

9) Usando as propriedades de esperança e variância, mostre que se $X\sim Binomial(n,p)$ entãro $\mathds{E}[X] = np$ e $Var(X) = np(1-p)$

\end{document}